\ifdefined\iscombined
	\problemset{Problem Set 2}
\else
\documentclass{article}
\usepackage{xparse}
\usepackage{tabularx}
\usepackage{changepage}
\usepackage{listings}
\usepackage{amsthm}
\usepackage{comment}
\usepackage{amsmath}
\usepackage{braket}
\usepackage{amsfonts}
\usepackage{sectsty}
\usepackage{hyperref}
\usepackage{fancyhdr}
\usepackage[top=1in,bottom=1in,left=1in,right=1in]{geometry}
\usepackage{float}
\usepackage{graphicx}
\usepackage{mathtools}

\date{
	\vspace{-.75em}
	{\large \today}
}

\title{
	\vspace*{-1.5em}
	{\Huge Problem Set 2} \\
	\vspace{-1em}
	\rule{\textwidth/2}{.01em} \\
	\vspace{-.75em}
}

\author{
	{\large Matthew Farah}
}

\hypersetup{
	colorlinks=true,
	linkcolor=blue,
	filecolor=magenta,
	urlcolor=blue,
	citecolor=blue,
	anchorcolor=blue
}

\fancyhead{}
\fancyhead[L]{\today}
\fancyhead[R]{Problem Set 2}

\setlength{\parindent}{0cm}

\newenvironment{solution}{
	\vspace{.5em}
	\par
	\begin{adjustwidth}{1.5em}{}
	\textbf{{Solution:}}
	\par
	\vspace{.5em}
}{
	\vspace{1em}
	\end{adjustwidth}
	\setcounter{step}{0}
}

\makeatother
\makeatletter

\newcounter{step}
\renewcommand{\thestep}{\Roman{step}}

\newenvironment{step}{
	\refstepcounter{step}
	\vspace{1em}\par\noindent
	\ignorespaces\noindent\textbf{(\thestep)}
	\ignorespaces
}{
	\par
	\vspace{1.5em}
}

\newcounter{problem}
\renewcommand{\theproblem}{\arabic{problem}}

\newenvironment{problem}[1][]{
	\newpage
	\refstepcounter{problem}
	\begin{adjustwidth}{1.5em}{}
	\noindent\textbf{Problem \theproblem: }\noindent\textit{#1}
	\par
	\phantomsection
	\addcontentsline{toc}{section}{\protect{\large{Problem \theproblem}}}
}{
	\vspace{.5em}
	\end{adjustwidth}
}


\newcounter{subproblem}
\renewcommand{\thesubproblem}{\alph{subproblem}}

\newenvironment{subproblem}{
	\setcounter{step}{0}
	\refstepcounter{subproblem}
	\vspace{1em}\par\noindent
	\phantomsection
	\addcontentsline{toc}{subsection}{\protect{\large{(\thesubproblem)}}}
	\begin{adjustwidth}{1.5em}{}
	\noindent{(\thesubproblem) }\ignorespaces
}{
	\par
	\vspace{1em}
	\end{adjustwidth}
}

%FIX: If I have time
\newcommand{\supp}[1][]{\text{supp}(#1)}

\sectionfont{\fontsize{12}{12}\bfseries}
\subsectionfont{\fontsize{10}{12}\normalfont}
\subsubsectionfont{\fontsize{10}{12}\normalfont}

\begin{document}

\maketitle
\pagestyle{fancy}

\newcolumntype{Y}{>{\centering\arraybackslash}X}

\tableofcontents

\fi


\begin{problem}[The centralizer of an element.]
	Let G be a group and let and let $a \in G$. The centralizer of $a$ is the set of all elements that commute with $a$:

	\begin{gather*}
		C(a) \coloneq \set{g \in G \colon ga = ag}
	\end{gather*}
\end{problem}

%FIX: (QUESTION 1.A)
% PROBABLY DONT NEED TO USE PROPOSITION 2.6.
%So ass

\begin{subproblem}
	In the group $S_5$, let $\sigma = (135)$. List all the elements of the centralizer $C(\sigma)$.
\end{subproblem}

\begin{solution}
	To find all elements in the centralizer of $\sigma$, we must find all elements, $\tau \in S_5$ such that $\sigma\tau = \tau\sigma$. \\

	Helping us in that endeavour is Proposition 2.6., which states that if $\tau$ satisfies $\supp[\sigma] \cap \supp[\tau] = \emptyset$, then $\tau$ and $\sigma$ commute. Clearly, $\sigma$ permutes 1, 3, and 5 while leaving 2 and 4 fixed. So, by definition of the support, $\supp[\sigma] = \set{1, 3, 5}$. Therefore, any $\tau$ satisfying $\supp[\tau] \subseteq \set{2, 4}$ must commute with $\sigma$. Evidently, there are only two such $\tau$ and so we can conclude that at least:

	\begin{gather*}
		\tau = (2, 4), \; \tau = () \\
	\end{gather*}

	must permute with $\sigma$ and so belong to the centralizer $C(\sigma)$.
\end{solution}

\begin{subproblem}
	Prove that $C(a) \le G$ for all $a \in G$.
\end{subproblem}

\begin{solution}
	First, let $a \in G$ be arbitrary. To prove that $C(a) \le G$, we must show that $C(a)$ is a subgroup of $G$, which we can do by the one-line subgroup test! \\

	The one line subgroup test states that if $xy^{-1} \in C(a)$ for all $x, y \in C(a)$, then $C(a)$ must be a subgroup of $G$. So, to that end, letting $x, y \in G$ be arbitrary, I will now proceed to show that $(xy^{-1})a = a(xy^{-1})$, hence meaning that $xy^{-1} \in C(a)$ and concluding the proof. \\

	By definition of $C(a)$, we know that $x, y \in a$ implies that $xa = ax$ and $ya = ay$. We also know, since $G$ is a group and $C(a)$ is clearly a subset of $G$, that $y$ must have an inverse in $G$. Letting $y^{-1} \in G$ be y's inverse, we can see that:

	\begin{align*}
		y^{-1}a &= y^{-1}(a(1)) &&\text{(By definition of the identity)} \\
		&= y^{-1}(a(yy^{-1}) &&\text{(By definition of inverses)} \\
		&= y^{-1}((ay)y^{-1}) &&\text{(By associativity)} \\
		&= (y^{-1}(ay))y^{-1} &&\text{(By associativity)} \\
		&= (y^{-1}(ya))y^{-1} &&\text{(Since $y \in C(a)$)} \\
		&= ((y^{-1}y)a)y^{-1} &&\text{(By associativity)} \\
		&= ((1)a)y^{-1}y^{-1} &&\text{(By definition of inverses)} \\
		&= ay^{-1} &&\text{(By definition of the identity)} \\
	\end{align*}

	Therefore, $y^{-1}a = ay^{-1}$, implying $y^{-1} \in C(a)$ and thereby making $xy^{-1} \in C(a)$ a corollary because:

	\begin{align*}
		(xy^{-1})a &= x(y^{-1}a) &&\text{(By associativity)} \\
		&= x(ay^{-1}) &&\text{(Since $y \in C(a)$)} \\
		&= (xa)y^{-1} &&\text{(By associativity)} \\
		&= (ax)y^{-1} &&\text{(Since $x \in C(a)$)} \\
		&= a(xy^{-1}) &&\text{(By associativity)} \\
	\end{align*}

	Immediately follows. Thus, by the one-line subgroup test, $C(a)$ is indeed a subgroup of $G$.
\end{solution}

\begin{subproblem}
	Two friends, Goofus and Gallant, are arguing about the following claim:

	\begin{gather*}
		\textit{``If } a, b \in G, \textit{then } C(a) \cap C(b) = C(ab)\text{''} \\
	\end{gather*}

	Goofus thinks ``$\subseteq$'' is true, but Gallant thinks ``$\supseteq$'' is true. Who is correct? Prove one inclusion, and give an explicit counterexample showing that the other inclusion is false in general.
\end{subproblem}

%FIX: QUESTION 1.C
% PROVIDE COUNTER EXAMPLE
% SOMETHING TO DO WITH GL
% ALSO SANITY CHECK LATER

\begin{solution}
	I contend, and will proceed to show, that $\subseteq$ is indeed the only correct inclusion. To do this, I intend to formally prove that $(C(a) \cap C(b)) \subseteq C(ab)$ for any arbitrary $a \in G$, for any arbitrary group $G$, and provide a specific group $G$ with elements $a, b \in G$ such that $(C(a) \cap C(b)) \supseteq C(ab)$ doesn't hold, making Gallant's inclusion false in general.

	\begin{step}
		Prove that $(C(a) \cap C(b)) \subseteq C(ab)$.
	\end{step}

	To prove this inclusion, I will show that $x \in (C(a) \cap C(b))$ implies that $x \in C(ab)$. So, let any such $x$ be arbitrary. \\

	By definition 

\end{solution}


\ifdefined\iscombined
	\newpage
\else
	\end{document}
\fi
